% ===========================================================================
%  sviluppatore/06-dashboard/04-pagina-nuova.tex
% ===========================================================================
\section{Aggiungere una pagina o una rotta}\label{sez:pagina-nuova}

\begin{pratica}{La procedura}
\begin{enumerate}
  \item Creare \file{dashboard/ui\_ext\_<lettera>.py}. La lettera decide
        l'ordine di caricamento; \file{c} e \file{d} sono libere.
  \item Esporre \simbolo{register(app, data\_layer, base\_path)} ---
        \emph{quella firma esatta}: è il contratto che il caricatore si
        aspetta.
  \item Dentro, registrare le rotte con \cod{@app.route}.
  \item Se serve una pagina, aggiungere il template in
        \file{dashboard/templates/} con \cod{\{\% extends "base.html" \%\}}.
  \item Se la rotta tocca l'account, metterla sotto uno dei prefissi di
        \cod{\_SERVE\_ACCOUNT} (\vediSez{permessi}).
  \item Riavviare. Non c'è nessun elenco da aggiornare.
\end{enumerate}
\end{pratica}

\subsection{Il modello da copiare}

\file{ui\_ext\_e.py} è l'esempio migliore, per due ragioni che valgono per
qualunque rotta nuova.

\subsubsection{Il docstring dichiara il contratto}

\begin{commento}{dashboard/ui\_ext\_e.py}
Rotte del gruppo E (Marketplace: vetrina, Sell yours, messaggi, feedback,
account) per la UI v2.

    GET  /api/e/sell-status   chi puo' pubblicare: piano (livello e nome),
                              amministratore, lucchetto del marketplace e, se
                              non si puo', il motivo nelle parole del sito.
    GET  /api/e/account-data  per la pagina My account, in una richiesta sola:
                              prodotti pubblicati, acquisti, richieste di
                              rimborso proprie e messaggi non letti.

Solo letture, e mai un errore verso la pagina: se il sito non risponde la
rotta dice «non so» (\cod{unknown: true}) e la pagina lo mostra come tale, invece
di far credere che il piano manchi.
\end{commento}

\subsubsection{«Non so» non è «no»}

È la regola di progetto più utile di tutto il capitolo, e vale ben oltre
questo file. Quando la dashboard interroga il sito e il sito non risponde,
la risposta sbagliata è \cod{can\_publish: false}: l'utente vede «non puoi
pubblicare», che è un'affermazione, e va a cercare un problema nel proprio
abbonamento. La risposta giusta è \cod{unknown: true}, che la pagina mostra
come incertezza.

Lo stesso schema ritorna in \percorsoHTTP{/api/e/account-data}, che riporta un campo
\cod{errors}: «se una lettura fallisce lo dice, così la pagina non scambia
\emph{non riuscito} per \emph{nessun acquisto}».

\begin{attenzione}{Un elenco vuoto è un'affermazione}
Restituire \cod{[]} quando non si è riusciti a leggere significa dire
all'utente «non hai acquisti». Per chi ne ha, è un errore che preoccupa, e
che porta ad aprire una contestazione. La distinzione fra «zero» e «non so»
va mantenuta fino alla pagina.
\end{attenzione}

\subsection{Una funzione di servizio che vale la pena copiare}

Le risposte JSON della dashboard contengono spesso date e \cod{Decimal} che
arrivano dal Postgres del sito, e che \cod{jsonify} non sa serializzare:

\begin{codiceInterno}{python}{Serializzazione difensiva}{dashboard/ui\_ext\_e.py:30}
def _plain(v):
    """Valori serializzabili in JSON (date in ISO, Decimal in float)."""
    if isinstance(v, (str, int, float, bool, type(None))):
        return v
    if isinstance(v, dict):
        return {str(k): _plain(x) for k, x in v.items()}
    if isinstance(v, (list, tuple, set)):
        return [_plain(x) for x in v]
    iso = getattr(v, "isoformat", None)
    if callable(iso):
        try:
            return iso()
        except Exception:  # noqa: BLE001
            pass
    try:
        return float(v)
    except Exception:  # noqa: BLE001
        return str(v)
\end{codiceInterno}

\noindent
L'ultimo ripiego è \cod{str(v)}: qualunque cosa arrivi, la pagina riceve
qualcosa invece di un \cod{500}. In una rotta di sola lettura è la scelta
giusta.

\subsection{Le due pagine speciali}

\begin{description}[leftmargin = 1.6em, style = nextline, itemsep = 5pt]
  \item[\file{ui\_ext\_b.py} fa un redirect in un \cod{before\_request}]
    Una richiesta a \percorsoHTTP{/ensemble\_custom?project=<id>} viene rediretta a
    \percorsoHTTP{/projects/<id>}. È il modo in cui i collegamenti vecchi continuano a
    funzionare dopo che una pagina ha cambiato indirizzo. Chi cerca perché un
    URL si trasforma in un altro lo trova lì, non nelle rotte.
  \item[\file{ui\_ext\_a.py} intercetta le installazioni in un \cod{after\_request}]
    Annota le installazioni fatte dallo Studio per i tutorial. Un
    \cod{after\_request} che osserva le risposte altrui è potente e
    invisibile: va dichiarato nel docstring del modulo, come lì è fatto.
\end{description}
